\documentclass[../main.tex]{subfiles}

\begin{document}

\section*{ABSTRACT}
\addcontentsline{toc}{section}{ABSTRACT}

This project presents the design and implementation of an integrated monitoring and control solution for modern poultry farming. The primary objective is to create a cost-effective, scalable, and automated system that enhances poultry health, optimizes resource usage, and ensures higher productivity through IoT-driven telemetry and adaptive control mechanisms. The hardware layer consists of distributed sensor modules that monitor critical environmental parameters such as temperature, humidity, air quality (ammonia/CO₂ levels), and light intensity within poultry houses. An Arduino Nano microcontroller node collects this telemetry and communicates with a central server via Wi-Fi , depending on deployment scale. The software pipeline incorporates a cloud-based data management system with real-time dashboards, alerts, and predictive analytics. Machine learning algorithms are applied to historical data to optimize feed schedules, predict health risks, and reduce mortality rates. A responsive web application developed using JavaScript provides farmers with interactive visualization, remote control of actuators, and anomaly detection alerts. Additionally, an adaptive control algorithm ensures that environmental conditions remain within optimal ranges to promote flock well-being. Initial phases of the project have demonstrated successful integration of IoT-based sensing, real-time telemetry collection, and basic actuator automation. Preliminary tests validate the system’s ability to regulate temperature and feeding cycles while providing reliable cloud-based monitoring. With low-cost hardware and open-source frameworks, this project proves to be a sustainable and accessible digital solution for both small-scale and commercial poultry farmers.

\textbf{Keywords}: \textit{Smart Farming, IoT, Telemetry, Adaptive Control, Poultry Management, Arduino Nano, Cloud Dashboard, Machine Learning, Precision Agriculture}



\end{document}
